\chapter{Conception de la base de données}
\label{chap:database}

\section{Introduction}
La base de données constitue le socle informationnel de la plateforme. Elle doit non seulement stocker les entités métier, mais surtout refléter avec précision les règles de gestion du parc informatique. Une conception relationnelle mal structurée aurait entraîné des doublons, des incohérences d'affectation, des états ambigus et une faible fiabilité des tableaux de bord. La phase de conception des données a donc été menée avec une attention particulière portée à la normalisation, à la cohérence métier et à la traçabilité.

\section{Modèle conceptuel de données}
Le modèle conceptuel de données a été construit à partir des objets métier centraux identifiés dans l'analyse : utilisateur, rôle, équipement, affectation, ticket, maintenance, notification et audit. La figure \ref{fig:mcd-core} présente une vue synthétique du noyau conceptuel.

\imagefigure{data-model-mcd.png}{MCD de la plateforme de gestion du parc informatique et des affectations}{fig:mcd-core}

Ce modèle montre que la donnée n'est pas organisée par écrans, mais par responsabilités métier. L'équipement devient le centre des flux opérationnels, tandis que l'affectation joue le rôle de pivot entre utilisateurs et actifs physiques.

\section{Modèle logique de données}
La traduction logique du MCD conduit à un ensemble d'entités relationnelles interconnectées. Dans l'implémentation Spring Boot observée, la couche \texttt{entity} comprend notamment :
\begin{itemize}
  \item \texttt{User}, \texttt{Role}, \texttt{Permission} ;
  \item \texttt{Equipment}, \texttt{Category}, \texttt{Supplier}, \texttt{Location}, \texttt{Department} ;
  \item \texttt{Assignment}, \texttt{EquipmentReturn}, \texttt{MaterialRequest} ;
  \item \texttt{Ticket}, \texttt{Maintenance} ;
  \item \texttt{Notification}, \texttt{AuditLog}, \texttt{ActivityLog} ;
  \item \texttt{RefreshToken}, \texttt{PasswordResetToken}, \texttt{EmailVérificationToken}.
\end{itemize}

Cette structuration logique montre que la base ne se limite pas à l'inventaire : elle embarque aussi la sécurité, le cycle de vie des sessions, les événements et la traçabilité.

\section{Schéma relationnel}
Le schéma relationnel peut être résumé par les associations majeures suivantes :
\begin{enumerate}
  \item un \texttt{User} possède un ou plusieurs \texttt{Role} ;
  \item un \texttt{Equipment} appartient à une \texttt{Category}, peut être localisé et affecté ;
  \item une \texttt{Assignment} relie un \texttt{User} à un \texttt{Equipment} avec un état métier ;
  \item une \texttt{MaterialRequest} formalise une demande de dotation avant affectation effective ;
  \item un \texttt{Ticket} peut concerner un équipement et ouvrir la voie à une \texttt{Maintenance} ;
  \item une \texttt{Notification} cible un utilisateur suite à un événement ;
  \item un \texttt{AuditLog} documente les actions sensibles sur les objets métier.
\end{enumerate}

Le tableau \ref{tab:schema-relationnel} formalise cette lecture relationnelle.

\begin{table}[H]
\centering
\caption{Synthèse du schéma relationnel}
\label{tab:schema-relationnel}
\begin{tabularx}{\textwidth}{>{\bfseries}p{3cm}X}
\toprule
Table & Relations dominantes \\
\midrule
users & liée à rôles, assignments, notifications, audit logs, tickets et maintenances \\
rôles & relation n:n avec users via table d'association \\
equipments & liée aux catégories, affectations, tickets, documents et statuts du parc \\
assignments & relie un bénéficiaire, un équipement et les acteurs de validation \\
material\_requests & relie un demandeur, un besoin de matériel, une priorité et une décision de traitement \\
tickets & rattaché les incidents à un équipement, un état et un niveau de priorité \\
maintenances & prolonge le ticket ou l'intervention planifiée sur un équipement \\
notifications & cible les utilisateurs selon le type d'événement métier \\
audit\_logs & historise les opérations critiques et les entités concernées \\
\bottomrule
\end{tabularx}
\end{table}

\section{Dictionnaire de données}
Le dictionnaire de données donne une vision plus fine des attributs essentiels. Le tableau \ref{tab:data-dictionary} présente un extrait représentatif des champs majeurs.

\begin{longtable}{p{2.6cm}p{3.2cm}p{2.5cm}p{6.2cm}}
\caption{Extrait du dictionnaire de données}\\
\toprule
\textbf{Entité} & \textbf{Attribut} & \textbf{Type} & \textbf{Description} \\
\midrule
\endfirsthead
\toprule
\textbf{Entité} & \textbf{Attribut} & \textbf{Type} & \textbf{Description} \\
\midrule
\endhead
User & id & UUID & identifiant technique unique \\
User & email & chaîne & identifiant de connexion unique \\
User & status & énumération & état du compte utilisateur \\
Role & name & chaîne & nom du rôle métier attribué \\
Equipment & inventoryCode & chaîne & code d'inventaire affiché dans l'application \\
Equipment & status & énumération & disponibilité actuelle du matériel \\
Equipment & warrantyEndDate & date & fin de garantie constructeur \\
Assignment & approved & booléen & indique si l'affectation a été validée \\
Assignment & expectedReturnDate & date & date théorique de restitution \\
MaterialRequest & materialType & chaîne & type de matériel demandé par l'employé bénéficiaire \\
MaterialRequest & status & enumeration & état de traitement: \texttt{PENDING}, \texttt{APPROVED}, \texttt{REJECTED}, \texttt{FULFILLED} ou \texttt{CANCELLED}. Dans les diagrammes, \texttt{FULFILLED} correspond à une demande clôturée. \\
Ticket & priority & énumération & niveau d'urgence du ticket \\
Ticket & status & énumération & état du traitement du ticket \\
Maintenance & scheduledAt & date/heure & date de planification de l'intervention \\
Maintenance & status & énumération & état d'avancement technique \\
Notification & title & chaîne & en-tête court du message utilisateur \\
AuditLog & description & texte & détail de l'action auditée \\
RefreshToken & expiresAt & date/heure & date d'expiration du jeton de renouvellement \\
\bottomrule
\label{tab:data-dictionary}
\end{longtable}

\section{Description des entités principales}
\subsection{Utilisateur, rôle et permissions}
Le triptyque \texttt{User} / \texttt{Role} / \texttt{Permission} permet d'implémenter un contrôle d'accès extensible. Au lieu de coder les autorisations directement dans l'interface, le modèle de données les matérialise explicitement. Cette approche favorise la maintenabilité et prépare l'extension vers de nouveaux rôles sans remise en cause de la structure globale.

\subsection{Équipement}
L'entité \texttt{Equipment} est le coeur du domaine métier. Elle concentre des attributs d'identification, de classification, de localisation, d'état et de garantie. Son rôle central justifie une attention particulière à la cohérence de ses statuts et à sa relation avec les affectations et les maintenances.

\subsection{Affectation}
L'entité \texttt{Assignment} formalise la mise à disposition d'un matériel à un bénéficiaire. Elle ne se contente pas de relier deux objets ; elle porte un véritable processus métier avec des dates, des acteurs de validation, un état et une logique de restitution. Elle matérialise donc le lien de gouvernance entre inventaire et usage.

\subsection{Demande de matériel}
L'entité \texttt{MaterialRequest} couvre l'étape amont de la dotation. Elle permet à un employé bénéficiaire d'exprimer un besoin avant qu'un responsable ne transforme éventuellement ce besoin en affectation. Elle porte le type de matériel demandé, le modèle souhaité, la justification, la priorité, le statut de traitement et l'acteur de revue.

\subsection{Ticket et maintenance}
Le couple \texttt{Ticket} / \texttt{Maintenance} permet de lier l'incident déclaré à l'action corrective ou préventive. Cette séparation est intéressante, car elle distingue l'événement fonctionnel du traitement technique, ce qui rend les rapports plus précis et les statistiques plus exploitables.

\section{Relations et cardinalités}
Les cardinalités essentielles peuvent être synthétisées comme suit :
\begin{itemize}
  \item un utilisateur peut avoir plusieurs rôles ;
  \item un équipement peut connaître plusieurs affectations successives mais une seule affectation active à un instant donné ;
  \item un utilisateur employé peut déposer plusieurs demandes de matériel dans le temps ;
  \item un équipement peut être associé à plusieurs tickets dans le temps ;
  \item un ticket peut aboutir à plusieurs opérations de maintenance selon la réalité terrain ;
  \item un utilisateur peut recevoir plusieurs notifications et générer plusieurs traces d'audit.
\end{itemize}

Le respect de ces cardinalités est un point critique de cohérence applicative. Il justifie le recours à une couche métier forte et à des contraintes d'intégrité adaptées.

\section{Contraintes de cohérence}
Plusieurs contraintes ont été prises en compte dans la conception :
\begin{enumerate}
  \item unicité de certains identifiants fonctionnels comme les adresses e-mail et les codes d'inventaire ;
  \item impossibilité logique d'avoir plusieurs affectations actives pour un même équipement ;
  \item obligation de traiter une demande de matériel par un acteur disposant des droits de gestion des affectations ;
  \item dépendance entre l'état d'un ticket et la possibilité d'ouvrir ou non une maintenance ;
  \item conservation d'un historique fiable des mouvements et validations ;
  \item sécurité des données liées aux sessions, jetons et vérifications d'e-mail.
\end{enumerate}

\section{Schémas PostgreSQL et scripts SQL}
Le projet embarque des scripts d'initialisation et de démonstration, notamment \texttt{01-init.sql} et \texttt{02-seed-demo-data.sql}. Ils complètent l'approche JPA/Hibernate en apportant une base de démarrage reproductible. Cette coexistence entre génération applicative et scripts d'environnement est utile dans un contexte de démonstration, de conteneurisation et de validation fonctionnelle.

\begin{codeblock}[H]
\caption{Script d'initialisation PostgreSQL}
\lstinputlisting[style=academiclst]{\fileInitSql}
\end{codeblock}

\section{Captures et exploitation de la base}
Dans le cadre de l'application, la base est sollicitée pour :
\begin{itemize}
  \item afficher le dashboard et les totaux métier ;
  \item filtrer les équipements disponibles ;
  \item restituer les listes d'utilisateurs et de maintenances ;
  \item générer les exports du module de rapports ;
  \item conserver la mémoire métier du système.
\end{itemize}

La qualité de la base de données influence donc directement la crédibilité de l'application, tant sur le plan fonctionnel que sur le plan analytique.

\section{Synthèse du chapitre}
La conception de la base de données reflète une lecture métier sérieuse du problème traité. Le schéma relationnel structure les équipements, les utilisateurs, les demandes de matériel, les affectations, la maintenance, les notifications et la sécurité dans un ensemble cohérent et extensible. Cette base solide conditionne la réussite de la réalisation logicielle détaillée dans le chapitre suivant.









